\documentclass[twocolumn]{aastex631}

\usepackage{amsmath}

\newcommand{\tic}{TIC~402710280}
\newcommand{\planet}{TIC~402710280.01}
\newcommand{\tess}{\textit{TESS}}
\newcommand{\gaia}{\textit{Gaia}}
\newcommand{\rearth}{R_\oplus}
\newcommand{\rsun}{R_\odot}
\newcommand{\msun}{M_\odot}

\shorttitle{A 36.34-day sub-Neptune candidate transiting TIC 402710280}
\shortauthors{Vora}

\begin{document}

\title{A 36.34-day Sub-Neptune Candidate Transiting the Nearby K Dwarf TIC 402710280}

\author[0009-0003-8778-9883]{Rishit Vora}
\affiliation{Independent Researcher}
\email{rishitv@protonmail.com}

\begin{abstract}
We report a transiting sub-Neptune candidate with a period of 36.34 days around \tic\ (TYC~9041-03677-1), a K dwarf of \tess\ magnitude 9.24 at 66.8~pc that the Transiting Exoplanet Survey Satellite (\tess) has observed only in full-frame images. Four transits, recorded between 2021 and 2026 in Sectors~39, 93, 100 and~102 and reduced by three independent pipelines, share a depth of $@@DEPTH@@$~ppm and a duration of $@@T14@@$~h. Two of them fall on consecutive cycles, and every shorter period predicts transits in covered data that are absent at $\geq @@ALIAS_MIN@@\sigma$; the period is $P = @@P@@$~d. A transit model gives $R_p/R_\star = @@K@@$ and a stellar density consistent with the host's. The candidate has a radius of $@@RP@@\,\rearth$ and an equilibrium temperature of $@@TEQ@@$~K. Difference-image centroids place the signal $1.5'' \pm 1.3''$ from the target. The only catalogued star close enough to contribute, a background star $3.5''$ away, would have to host an eclipsing binary whose eclipse shape is disfavoured for any main-sequence pair ($\Delta\chi^2 \geq 142$). No secondary eclipse is detected, and \texttt{TRICERATOPS} gives a false-positive probability of $0.003 \pm 0.004$. No high-resolution imaging exists, so companions within about $1''$ are constrained only by \gaia\ astrometry. The mid-times scatter about a linear ephemeris by @@RMS@@~min (rms). Quick-Look Pipeline light curves reproduce the 29-min excess in the interval between the two 2026 transits, so it does not arise in our photometry; we report the scatter without choosing between an instrumental and a dynamical origin. The light curve also contains a $@@SS_DEPTH@@$~ppm signal with a period of 3.78~days whose source, the target or the background star, the present data cannot identify. We predict transit times, including one in \tess\ Sector~114 on 2027 February~19.
\end{abstract}

\keywords{Exoplanets (498) --- Transit photometry (1709) --- Mini Neptunes (1063) --- K dwarf stars (876)}

\section{Introduction} \label{sec:intro}

Sub-Neptunes with orbital periods longer than a month around bright K dwarfs are rare among known transiting planets, yet they are the temperate members of the population for which both precise masses and atmospheric spectra are within reach: the stars are bright enough for high-precision radial velocities, and their small radii deepen both the transit and its atmospheric signal. A single sector of the Transiting Exoplanet Survey Satellite \citep[\tess;][]{Ricker2015}, about 27~d long, records at most one transit of such a planet. Finding them therefore means combining sectors observed years apart, which became possible only as the full-frame images accumulated.

Stars observed only in full-frame images lie outside the Science Processing Operations Center (SPOC) search of two-minute data \citep{Jenkins2016}. Their long-period planets must be found in full-frame light curves, such as those of the Quick-Look Pipeline \citep[QLP;][]{Huang2020a, Huang2020b} and TESS-SPOC \citep{Caldwell2020}, and in photometry of the most recent images, which these products cover only weeks to months later. The star studied here illustrates the difficulty: until 2026 its public light curves contained two transits four years apart, and the one predicted between them fell on a scattered-light ramp, so the period could not be measured from public data.

Here we report a sub-Neptune candidate with a period of 36.34~days around \tic, a K dwarf at 67~pc. Its first two transits were identified in public full-frame light curves; the period was determined from our photometry of the 2026 full-frame images. As of 2026 October~14, the star is neither a \tess\ Object of Interest \citep[TOI;][]{Guerrero2021} nor a community TOI, and no planet parameters, follow-up observations or literature references are associated with it in the Exoplanet Follow-up Observing Program (ExoFOP) archive, the NASA Exoplanet Archive \citep{Akeson2013}, or the NASA Astrophysics Data System. We refer to the candidate as \planet, the designation ExoFOP assigns to the first community TOI on a star; we will submit it as a community TOI upon publication in a refereed journal, as ExoFOP requires for new candidates.

\section{Host Star} \label{sec:star}

Table~\ref{tab:star} lists the properties of \tic. We adopt the stellar parameters of the \tess\ Input Catalog v8.2 \citep[TIC;][]{Stassun2019}. The \gaia\ DR3 \citep{Gaia2016, GaiaDR3} astrophysical parameters agree with them: the GSP-Phot effective temperature is 5089~K \citep{Andrae2023}, and the FLAME radius and mass are $0.843 \pm 0.018\,\rsun$ and $0.81 \pm 0.04\,\msun$ \citep{Creevey2023}, within $1\sigma$ of the TIC values. Adopting the FLAME radius would increase the radius of the candidate by 3\%.

\gaia\ DR3 reports a renormalised unit weight error \citep[RUWE;][]{Lindegren2021} of $1.00$, no multi-peaked image windows (\texttt{ipd\_frac\_multi\_peak} $= 0$), and no non-single-star solution. These indicate a single, well-behaved astrometric source, although they do not exclude a companion at small separation and low mass ratio.

The light curves show no coherent rotational modulation. The strongest Lomb--Scargle peaks of the individual sectors lie between 5 and 13~d, differ from sector to sector, and do not exceed 0.1\% peak to peak in the TESS-SPOC light curves, whose systematics correction suppresses longer periods. No light curve contains two or more consecutive cadences more than $4\sigma$ above its running median, except during one scattered-light episode in Sector~93 (BTJD~3830.6, see \S\ref{sec:data}), where the excess follows the background; we find no flares. No stellar signal is therefore available with which the transits could be phased.

\begin{deluxetable}{lcc}
\tablecaption{Properties of \tic \label{tab:star}}
\tablehead{\colhead{Parameter} & \colhead{Value} & \colhead{Source}}
\startdata
R.A. (J2000) & 16:34:05.67 & TIC \\
Decl. (J2000) & $-$63:32:58.8 & TIC \\
\gaia\ DR3 ID & 5828441289812582400 & \gaia \\
$\mu_\alpha \cos\delta$ (mas yr$^{-1}$) & $59.909$ & \gaia \\
$\mu_\delta$ (mas yr$^{-1}$) & $13.685$ & \gaia \\
Parallax (mas) & $14.964 \pm 0.014$ & \gaia \\
Distance (pc) & $66.83 \pm 0.06$ & \gaia \\
RUWE & $1.00$ & \gaia \\
$T$ (mag) & $9.243 \pm 0.006$ & TIC \\
$G$ (mag) & $9.781$ & \gaia \\
$G_{\rm BP} - G_{\rm RP}$ (mag) & $1.023$ & \gaia \\
$V$ (mag) & $9.982 \pm 0.004$ & TIC \\
$K_s$ (mag) & $8.026$ & TIC \\
Contamination ratio & $0.060$ & TIC \\
$T_{\rm eff}$ (K) & $5182 \pm 121$ & TIC \\
$\log g$ (cgs) & $4.56 \pm 0.08$ & TIC \\
$R_\star$ ($\rsun$) & $0.817 \pm 0.047$ & TIC \\
$M_\star$ ($\msun$) & $0.88 \pm 0.11$ & TIC \\
$\rho_\star$ ($\rho_\odot$) & $1.61 \pm 0.38$ & derived \\
$L_\star$ ($L_\odot$) & $0.430 \pm 0.005$ & \gaia\ FLAME \\
\enddata
\tablecomments{$\rho_\star$ is computed from the TIC mass and radius. The contamination ratio is the TIC estimate of the flux from other stars relative to the target's within a \tess\ aperture.}
\end{deluxetable}

\section{Observations and Photometry} \label{sec:data}

\tess\ observed \tic\ in Sectors~12 (2019), 39 (2021), 66 (2023), 93 (2025) and 100--103 (2026), always in full-frame images, at cadences of 1800~s (prime mission), 600~s (first extended mission) and 200~s thereafter. We use the TESS-SPOC light curves \citep{Caldwell2020, TESSSPOC} for Sectors~12, 39 and~66 (pre-search data conditioning flux) and the QLP light curve \citep{Huang2020a, QLPdata} for Sector~93 (systematics-corrected flux). We extracted the light curves of Sectors~100--103 from the SPOC-calibrated full-frame images, using $21 \times 21$-pixel cutouts \citep{Brasseur2019}. The background of each cadence is the median of the faintest 40\% of pixels more than 4~pixels from the target. The aperture is a circle whose radius, between 1 and 4~pixels in steps of 0.5, minimises the one-hour scatter of the detrended light curve; the adopted radii are 1.5--2.5~pixels. Cadences flagged by SPOC are removed, as are scattered-light glints, in which the background exceeds its one-day running median by more than eight times its median absolute deviation. The first version of this paper stated that no pipeline had processed these sectors; in fact, QLP light curves of them were released at MAST between 2026 March and July, but they are absent from the public cloud copy of the archive that our pipeline reads. We use them only to check our timing (\S\ref{sec:ttv}).

The time stamps of the calibrated full-frame images are barycentric for a reference position on the detector rather than for the target. We recomputed them for the target's position, a correction of $-48$~s and $-20$~s at the Sector~100 and 102 transits. The TESS-SPOC and QLP light curves we fit are corrected for the target already: their time stamps precede those of the same cadences in the calibrated images by 20.7~s, in agreement with the 20--23~s we compute.

The TESS-SPOC fluxes are corrected for the flux of neighbouring stars (\texttt{CROWDSAP} $= 0.92$--$0.94$); we apply no such correction to our photometry. The depths measured in the two kinds of light curve agree within their 4--7\% uncertainties (\S\ref{sec:times}), so any difference in dilution between the reductions is smaller than the measurement precision.

For transit analysis, each event is extracted in a window of $\pm 0.8$~d and normalised by a quadratic function of time fitted to the out-of-transit data.

\section{Detection and Period} \label{sec:period}

The candidate was found in a search of 10\,324\,827 public light curves (QLP light curves of Sectors~1--98 and SPOC two-minute light curves of nine sectors) for individual transit-like events. Box filters of ten durations between 0.75 and 18~h were applied, with a detection threshold of $6\sigma$ relative to the scatter of each filtered series. After removal of systematics shared by many stars, events at data edges and gaps, and non-planetary shapes, 2263 stars retained single transit-like events. The 142 of these brighter than $T = 10.5$ were vetted individually at the pixel level and against catalogues of known objects, and the 27 that survived without a known period were searched for further transits in our photometry of Sectors~99--108.

For \tic\ the public data contain two dips of matching depth and duration, in Sectors~39 and 93, separated by 1453.58~d. Our photometry revealed two further transits, in Sectors~100 and 102, separated by 36.36~d (Figure~\ref{fig:transits}). Because these two transits are consecutive, the period cannot be a multiple of 36.36~d; the remaining possibilities are its integer fractions. We tested $P/2$ through $P/5$ by measuring the flux at every transit time they predict within the data, excluding the four observed events. $P/2$ predicts transits at @@N_P2@@ covered epochs (Sectors~@@P2_SECTORS@@), each absent at $\geq @@P2_MIN@@\sigma$; $P/3$ predicts @@N_P3@@, absent at $\geq @@P3_MIN@@\sigma$; and $P/4$ and $P/5$ predict @@N_P45@@ more, of which the least significant exclusion is $@@ALIAS_MIN@@\sigma$. The adopted period predicts no transit in Sector~12, whose data (BTJD 1624.9--1652.9, where BTJD $=$ BJD$_{\rm TDB} - 2457000$) fall between predicted transits at BTJD~1621.5 and 1657.8, and none in Sector~101; its Sector~103 transit, at BTJD~4165.26, falls in the data gap between BTJD~4164.44 and 4166.79. A period of 36.34~d is therefore the only one consistent with the covered epochs.

The transit predicted in Sector~66 (BTJD~3111.40) coincides with an interval of strong scattered light (Figure~\ref{fig:s66}). The background falls from $1.5 \times 10^4$ to $2.9 \times 10^3$~e$^-$~s$^{-1}$ across the window, quality-flagged cadences end @@S66_FIRST_H@@~h before the predicted mid-time (@@S66_PRE_INGRESS_H@@~h before ingress), and QLP flags every cadence within 12~h of it. The transit is therefore covered, but the baseline before it is not. Fitted with the shape of \S\ref{sec:fit} and a free mid-time, the event has depths of @@S66_DEPTHS@@~ppm (uncertainties @@S66_DERR@@~ppm) for constant, linear and quadratic baselines, and falls @@S66_OC_RANGE@@~min from the linear ephemeris of \S\ref{sec:ttv}. Its timing is stable, but its depth depends on the one-sided baseline. We therefore exclude Sector~66 from all fits and from the ephemeris.

\begin{figure*}
\epsscale{1.15}
\plotone{figures/fig1.pdf}
\caption{Transits of \planet. Grey points are individual cadences; blue points are 20-minute bins, with error bars equal to the standard error of the mean of the cadences in each bin. Each window is centred on its measured mid-time (Table~\ref{tab:times}), and the red curve is the median transit model of \S\ref{sec:fit}; the lower panels show the binned residuals. The per-cadence scatter of our Sector~100 and 102 photometry (Table~\ref{tab:times}) is larger than that of the QLP light curve at the same cadence because our apertures are smaller and are not optimised per cadence. The right-hand panel shows the Sector~66 window at the time predicted by the linear ephemeris; data with zero quality flags begin at the dotted line, and the window is not used in any fit (\S\ref{sec:period}).\label{fig:transits}}
\end{figure*}

\begin{figure}
\plotone{figures/fig2.pdf}
\caption{TESS-SPOC data of Sector~66 around the predicted transit, whose duration is shaded and whose ingress and egress are marked by dashed lines. From top: simple-aperture flux, with cadences flagged for scattered light in grey; pre-search data conditioning flux, quality 0 only; and the background.\label{fig:s66}}
\end{figure}

\section{Transit Analysis} \label{sec:analysis}

\subsection{Individual Transit Times and Depths} \label{sec:times}

Each transit is timed by fixing the shape to the median model of \S\ref{sec:fit} and scanning the mid-time on a 0.5-minute grid over $\pm 2$~h, refitting a quadratic baseline over the $\pm 0.8$~d window at each step. The $1\sigma$ uncertainty is the half-width of the $\Delta\chi^2 = 1$ interval multiplied by the red-noise factor $\beta$ \citep{Pont2006, Winn2008}, the median ratio of the binned residual scatter (10--60-minute bins) to its white-noise expectation. The depth of each event is measured with the shape fixed and a free amplitude. Table~\ref{tab:times} lists the results.

The four depths agree ($\chi^2 = @@DEPTH_CHI2@@$ for 3 degrees of freedom) although they come from three independent reductions with different apertures, cadences and crowding treatments. The agreement argues against an origin in any one pipeline and against aperture-dependent dilution. Their weighted mean is $@@DEPTH@@$~ppm.

\begin{deluxetable*}{lcccccccc}
\tablecaption{Individual transits of \planet \label{tab:times}}
\tablehead{\colhead{Sector} & \colhead{Light curve} & \colhead{Cadence (s)} & \colhead{$\sigma$ (ppm)} & \colhead{$\beta$} & \colhead{Epoch} & \colhead{$T_c$ (BJD$_{\rm TDB}$)} & \colhead{$O - C$ (min)} & \colhead{Depth (ppm)}}
\startdata
@@TIMES_ROWS@@
\enddata
\tablecomments{$\sigma$ is the out-of-transit scatter per cadence and $\beta$ the red-noise factor of each window; every window spans $\pm 0.8$~d and has a quadratic baseline. $O-C$ is relative to the linear ephemeris of Table~\ref{tab:params}. Depths are the model depth at mid-transit for a free amplitude with the shape fixed.}
\end{deluxetable*}

\subsection{Transit Shape} \label{sec:fit}

We model the four transits jointly with \texttt{batman} \citep{Kreidberg2015}, integrating over each light curve's exposure time and centring each window on its measured mid-time, so that timing scatter (\S\ref{sec:ttv}) does not smear the stacked profile. The free parameters are $R_p/R_\star$, $b$ and $\log_{10}\rho_\star$, with uniform priors; $a/R_\star$ follows from $\rho_\star$ and $P$ \citep{Seager2003} for a circular orbit. Quadratic limb-darkening coefficients are fixed at $(u_1, u_2) = (0.39, 0.21)$, interpolated in $T_{\rm eff}$ from the \tess-band tables of \citet{Claret2017} at $\log g = 4.5$. Each window's flux uncertainty is its out-of-transit point-to-point scatter multiplied by its $\beta$. We sample the posterior with \texttt{emcee} \citep{Goodman2010, ForemanMackey2013}, using 40 walkers for 30\,000 steps and discarding the first third; the chain spans more than @@NTAU@@ integrated autocorrelation times for every parameter. Table~\ref{tab:params} gives the medians and 68\% intervals. The best-fitting model has a reduced $\chi^2$ of @@REDCHI2@@ (@@NPTS@@ data points), so the $\beta$-scaled uncertainties are slightly conservative.

The impact parameter is below @@B95@@ at 95\% confidence. The fitted stellar density, $@@RHO@@\,\rho_\odot$, agrees with the value of $1.61 \pm 0.38\,\rho_\odot$ derived from the TIC mass and radius; the transit is therefore consistent with a circular orbit around a star of the catalogued density. Fitting the limb darkening as well, with the uninformative parametrisation of \citet{Kipping2013}, changes $R_p/R_\star$ to @@K_FREELD@@ and $b$ to @@B_FREELD@@, both within $1\sigma$ of the adopted values.

\subsection{Ephemeris and Timing Residuals} \label{sec:ttv}

A weighted linear fit to the four mid-times gives $P = @@P@@$~d and $T_0 = @@T0@@$~BJD$_{\rm TDB}$, with the reference epoch at the Sector~93 transit, near the weighted mean epoch. The residuals (Table~\ref{tab:times}, Figure~\ref{fig:oc}) are @@OC_LIST@@~min, with an rms of @@RMS@@~min, much larger than the @@ERR_RANGE@@-min formal uncertainties ($\chi^2 = @@EPH_CHI2@@$ for 2 degrees of freedom). We scale the ephemeris uncertainties by $\sqrt{\chi^2/\nu}$. The largest single contribution is the interval between the consecutive Sector~100 and 102 transits, which exceeds $P$ by @@INTERVAL_EXCESS@@~min.

We searched for an instrumental cause. Varying the baseline (constant, linear or quadratic) and the window half-width (0.6, 0.8 or 1.0~d) moves no mid-time by more than 2.3~min. The Sector~93 time is the same, within 1.9~min, in every QLP flux series (the systematics-corrected flux, the raw flux, and the detrended fluxes of the optimal, small and large apertures), and that window is complete (688 of 691 cadences, no gap longer than 7~min), with no background or centroid feature at the transit. The time stamps of all four light curves refer to the target's barycentric position (\S\ref{sec:data}). Masking the transits of the 3.78-day signal (\S\ref{sec:second}), one of which falls 4~h before the Sector~93 transit, changes no mid-time by more than 0.1~min.

The QLP light curves of Sectors~100 and~102 test the two mid-times from our photometry with an independent reduction of the same images. Their time stamps need a correction first. Matched to the full-frame images by cadence number, they precede the mid-exposure times of the images, corrected to the target, by @@QLP_OFF100@@ and @@QLP_OFF102@@~s at the two transits; in Sector~93 they agree within @@QLP_OFF93@@~s. Corrected and timed as in \S\ref{sec:times}, the QLP transits fall @@QLP_EARLY100@@ and @@QLP_EARLY102@@~min earlier than ours (@@QLP_SIG100@@ and @@QLP_SIG102@@ times the combined uncertainty), and the interval between them exceeds $P$ by @@QLP_INTERVAL@@~min, against @@INTERVAL_EXCESS@@~min in our photometry. Our photometry of the Sector~39 and 93 transits, reduced in the same way, gives mid-times within @@OWN_MAX@@~min of the TESS-SPOC and QLP values. The excess in the Sector~100--102 interval therefore does not arise in our reductions. The QLP depths, $@@QLP_DEP100@@$ and $@@QLP_DEP102@@$~ppm, are lower than ours by @@QLP_DSIG100@@ and @@QLP_DSIG102@@ times the combined uncertainty; we have not identified the cause, and neither QLP light curve enters the fits. We therefore identify no instrumental cause, and four epochs cannot distinguish timing systematics from transit-timing variations induced by an additional body. Additional transits will show whether the residuals are periodic. The predictions in \S\ref{sec:predictions} include the observed scatter.

\begin{figure}
\plotone{figures/fig3.pdf}
\caption{Transit-timing residuals relative to the linear ephemeris. Symbols denote the light-curve source: TESS-SPOC (square), QLP (diamond) and our photometry (open circles); the formal uncertainties, @@ERR_RANGE@@~min, are smaller than the symbols. The shaded band is the $1\sigma$ uncertainty of the ephemeris; the dotted curves add the @@RMS@@-min rms scatter of the residuals in quadrature. Dashed lines mark epochs 57 (2027 February~19, \tess\ Sector~114) and 59 (2027 May~3, observable from La Silla).\label{fig:oc}}
\end{figure}

\begin{deluxetable}{lc}
\tablecaption{Parameters of \planet \label{tab:params}}
\tablehead{\colhead{Parameter} & \colhead{Value}}
\startdata
\multicolumn{2}{c}{\textit{Ephemeris}} \\
$P$ (d) & $@@P@@$ \\
$T_0$ (BJD$_{\rm TDB}$) & $@@T0@@$ \\
Timing residual rms (min) & @@RMS@@ \\
\multicolumn{2}{c}{\textit{Fitted}} \\
$R_p/R_\star$ & $@@K_FULL@@$ \\
$b$ & $@@B@@$ \\
$\rho_\star$ ($\rho_\odot$) & $@@RHO@@$ \\
\multicolumn{2}{c}{\textit{Fixed}} \\
$e$ & 0 \\
$(u_1, u_2)$ & $(0.39, 0.21)$ \\
\multicolumn{2}{c}{\textit{Derived}} \\
Depth at mid-transit (ppm) & $@@DEPTH@@$ \\
$T_{14}$ (h) & $@@T14@@$ \\
$a/R_\star$ & $@@ARS@@$ \\
$R_p$ ($\rearth$) & $@@RP@@$ \\
$a$ (AU) & $@@A_AU@@$ \\
$S$ ($S_\oplus$) & $@@S@@$ \\
$T_{\rm eq}$ (K) & $@@TEQ@@$ \\
\enddata
\tablecomments{Medians with 68\% intervals. The depth is the weighted mean of the four measurements of Table~\ref{tab:times}. $T_{\rm eq}$ assumes zero Bond albedo and full heat redistribution. $R_p$, $a$, $S$ and $T_{\rm eq}$ use the stellar parameters of Table~\ref{tab:star}.}
\end{deluxetable}

\section{Vetting} \label{sec:vetting}

Without high-resolution imaging or spectroscopy, the light curves and the \gaia\ catalogue can test four things: whether the signal comes from the target's position (\S\ref{sec:centroid}), whether the one catalogued star near enough to contribute can produce it (\S\ref{sec:blend}), whether the target is itself an eclipsing binary (\S\ref{sec:secondary}), and how probable the remaining false-positive scenarios are given the field (\S\ref{sec:fpp}). They cannot exclude an unresolved companion closer than about $1''$.

\subsection{Location of the Signal} \label{sec:centroid}

For each transit we form a difference image (out-of-transit minus in-transit) from the full-frame image cutouts and fit a Gaussian point-spread function to it \citep[cf.][]{Bryson2013, Twicken2018}. Each offset is measured relative to the target's own fitted position in the out-of-transit image, which removes errors in the image's world-coordinate solution and the effect of proper motion. Combining the four transits gives an offset of $1.5'' \pm 1.3''$ from \tic\ (Figure~\ref{fig:diff}). We validated the method on 40 confirmed \tess\ planets: it placed 30 on their host star, left 5 unresolved from a nearby source, returned an offset without a catalogued counterpart for 3, failed for 2, and placed none on a neighbouring star.

The centroid is $1.2\sigma$ from the target and $1.6\sigma$ from one neighbour: it excludes every other \gaia\ DR3 source bright enough to produce the observed depth, each of which lies $21\sigma$ or more away, but not this one. The neighbour, \gaia\ DR3 5828441289812582656 ($G = 16.15$, $\Delta G = 6.37$), has a parallax of $0.69 \pm 0.08$~mas, which places it in the background. Owing to the target's proper motion, its separation increases from $3.10''$ at the \gaia\ epoch (2016.0) to $3.4''$--$3.6''$ during the transits.

\begin{figure*}
\epsscale{1.1}
\plotone{figures/fig5.pdf}
\caption{Location of the signal. Left and centre: the out-of-transit image and the difference image (out of transit minus in transit) of the Sector~39 transit, the one with the highest signal-to-noise ratio; the plus marks \tic\ and the cross the background neighbour, which are not separated at the 21$''$ pixel scale. Right: the centroid offsets of the four transits (grey) and their combination (blue), with $1\sigma$ (solid) and $3\sigma$ (dotted) circles, relative to \tic; the cross marks the neighbour at the mean epoch of the transits.\label{fig:diff}}
\end{figure*}

\subsection{Blended Eclipsing Binary} \label{sec:blend}

Diluted by \tic, the neighbour could reproduce the observed depth only through an eclipse of about 74\% of its own light. We fitted the four transits with eclipse models of the neighbour for companion-to-star radius ratios $k$ from 0.3 to 10, with the impact parameter and $a/R_\star$ free, the dilution free but bounded at 0.4\% of the total flux (40\% above the star's share, to allow for its unknown \tess-band colour), and three limb-darkening laws (Figure~\ref{fig:blend}). For $k \leq 1.2$ every eclipse model is worse than the planet model by $\Delta\chi^2 \geq 142$ for 2227 data points. Larger companions fit better ($\Delta\chi^2 = 113$, 70, 44, 30 and 18 for $k = 1.5$, 2, 3, 5 and 10), because a total eclipse with a short ingress resembles a transit, but only if the eclipsed star supplies about 70\% of the neighbour's light while a companion 1.5 to 10 times its radius supplies the rest. The companion's surface brightness would then have to be below about 20\% ($k = 1.5$) to 0.5\% ($k = 10$) of the eclipsed star's, which no main-sequence pair satisfies; only a hot, compact eclipsed star could. A blend with the neighbour is therefore disfavoured, and a ground-based light curve with an aperture that excludes it would settle the question directly.

\begin{figure}
\plotone{figures/fig4.pdf}
\caption{Diluted-eclipse models for the background neighbour, compared with the planet model. For each companion radius ratio $k$, the curve gives the smallest $\chi^2$ difference over impact parameter, $a/R_\star$, dilution and limb darkening, and the band the range over the three limb-darkening laws. The shaded region, $k \leq 1.2$, contains every main-sequence pair.\label{fig:blend}}
\end{figure}

\subsection{Secondary Eclipse and Odd--Even Depths} \label{sec:secondary}

The light curve at orbital phase 0.5, sampled in Sectors~39, 101 and~102, shows no eclipse: the weighted depth is $@@SEC_DEPTH@@$~ppm, below @@SEC_3SIG@@~ppm at $3\sigma$. A search over phases 0.3 to 0.7, which covers secondary eclipses of eccentric orbits with $|e\cos\omega| \lesssim 0.3$, finds no consistent signal: the 3$\sigma$ upper limit is below 160~ppm at half of the phases and below 365~ppm at all of them. The most significant box, at phase 0.449 ($204 \pm 53$~ppm), is dominated by data within 0.5~d of the end of Sector~93 and is not reproduced in Sectors~39 and 101. These limits exclude an eclipsing companion of comparable surface brightness to the target. Only one of the four transits (Sector~100) falls on an odd epoch, so a comparison of odd and even depths is not informative; the agreement of all four depths (\S\ref{sec:times}) nonetheless disfavours an eclipsing binary at twice the period with unequal eclipses.

\subsection{Statistical False-Positive Probability} \label{sec:fpp}

Statistical validation with \texttt{TRICERATOPS} \citep{Giacalone2020, Giacalone2021} requires high-resolution imaging to constrain unresolved companions, and none exists for this star; the probabilities below are therefore upper bounds on the evidence and do not validate the candidate. We computed scenario probabilities using the four transits stacked on their measured mid-times (10-minute bins), a depth of 2100~ppm, the photometric aperture of each sector, and the \gaia\ DR3 field. Ten independent runs give a false-positive probability ${\rm FPP} = 0.003 \pm 0.004$ (largest value 0.012) and a nearby false-positive probability ${\rm NFPP} < 10^{-6}$.\footnote{With the photometric apertures NFPP is numerically zero ($<10^{-20}$); with \texttt{TRICERATOPS}'s default $5 \times 5$-pixel apertures, ${\rm FPP} = 0.001$ and ${\rm NFPP} = 6 \times 10^{-7}$.} The planet hypothesis on the target carries 85\% of the probability and a planet diluted by an unresolved bound companion 14\%; the leading false-positive scenario, a planet orbiting an unresolved companion, carries 0.3\%. These values satisfy the numerical criteria of \citet{Giacalone2021}, ${\rm FPP} < 0.015$ and ${\rm NFPP} < 0.001$, but without imaging we classify \planet\ as a planet candidate.

\section{A Second Periodic Signal} \label{sec:second}

With the transits of \planet\ masked, a box-least-squares search \citep{Kovacs2002} of the light curve finds a second periodic signal, with $P = @@SS_P@@$~d, a depth of $@@SS_DEPTH@@$~ppm and a duration of about @@SS_T14@@~h (Figure~\ref{fig:second}). It is detected independently in Sectors~12--93, as the strongest periodicity (signal detection efficiency 14.4), and in Sectors~100--103, as the second strongest after a period five times longer; the two ephemerides agree in phase across 650 cycles. The depths in the eight sectors agree with one another ($\chi^2 = @@SS_SECTOR_CHI2@@$ for 7 degrees of freedom; Figure~\ref{fig:second}), the odd and even transits agree ($@@SS_ODD@@$ and $@@SS_EVEN@@$~ppm), and no eclipse is seen at phase 0.5 ($@@SS_SEC@@$~ppm).

A transit model with free stellar density gives @@SS_FIT@@. @@SS_INTERP@@ The difference-image centroid of this weaker signal ($5.2'' \pm 6.7''$ from the target) cannot separate the target from the background neighbour, on which it would require an eclipse of about @@SS_NEIGH@@\% of the neighbour's light. @@SS_FPP@@ We therefore do not claim this signal as a planet candidate. A ground-based light curve that resolves the neighbour would show a @@SS_NEIGH@@\% eclipse on it if it is the source.

\begin{figure}
\plotone{figures/fig6.pdf}
\caption{The 3.78-day signal. Top: all eight sectors folded on its ephemeris, with the transits of \planet\ masked, in 15-minute bins (error bars are standard errors of the mean), and the median transit model. Bottom: the box depth in each sector, with uncertainties from boxes placed at 400 random phases; the line is the mean over all sectors.\label{fig:second}}
\end{figure}

\section{Discussion} \label{sec:discussion}

With $R_p = @@RP@@\,\rearth$, \planet\ lies above the radius valley \citep{Fulton2017}, among sub-Neptunes that probably retain a substantial volatile envelope. At $@@A_AU_SHORT@@$~AU it receives ten times Earth's insolation. The mass--radius relation of \citet{ChenKipping2017} predicts $M_p \approx 14\,M_\oplus$ and a radial-velocity semi-amplitude of $K \approx 2.8$~m~s$^{-1}$, within reach of current high-precision spectrographs for a $V = 10.0$ star.

Three observations would most efficiently test the candidate: high-resolution imaging, to exclude unresolved companions and complete the statistical validation; seeing-limited photometry of a transit with an aperture that excludes the $3.5''$ neighbour, which would locate both signals; and radial velocities, to measure the mass. Each could overturn the present interpretation. A companion within $1''$ bright enough to dilute the transit would change the radius and the false-positive probability; a 74\% eclipse of the neighbour at the time of a transit would identify the signal as a blend; and a semi-amplitude far above the few metres per second expected for a sub-Neptune would indicate a stellar or substellar companion.

\section{Predicted Transits} \label{sec:predictions}

Table~\ref{tab:predictions} lists three of the next transits predicted by the ephemeris of Table~\ref{tab:params}; the next 30 are listed in the data record \citep{DataRecord}. The quoted uncertainty, $\sigma_{\rm pred} = (\sigma_{T_0}^2 + n^2\sigma_P^2 + \sigma_{\rm rms}^2)^{1/2}$, adds the observed residual scatter $\sigma_{\rm rms} = @@RMS@@$~min to the ephemeris uncertainty, so that it describes an individual transit rather than the mean ephemeris. \tess\ is expected to observe \tic\ in Sectors~113 and~114 \citep{Burke2020}; the transit of 2027 February~19 is predicted to fall in Sector~114. The transit of 2027 May~3 is observable in full from the southern hemisphere. The test of these predictions, with the data to be used and the rules that decide its outcome, was registered before the data existed \citep{PreReg}.

\begin{deluxetable*}{ccccl}
\tablecaption{Predicted transits of \planet \label{tab:predictions}}
\tablehead{\colhead{Epoch} & \colhead{$T_c$ (BJD$_{\rm TDB}$)} & \colhead{$\sigma_{\rm pred}$ (min)} & \colhead{UTC at La Silla} & \colhead{Note}}
\startdata
@@PRED_ROWS@@
\enddata
\tablecomments{Epochs count cycles from the Sector~39 transit. UTC times are topocentric at La Silla Observatory: the light-travel time and TDB$-$UTC are removed.}
\end{deluxetable*}

A transit absent from the predicted window, in data that cover it, would falsify the reported period.

\section*{Acknowledgments}
This paper includes data collected by the \tess\ mission, which are publicly available from the Mikulski Archive for Space Telescopes (MAST). Funding for the \tess\ mission is provided by NASA's Science Mission Directorate. This work has made use of data from the European Space Agency (ESA) mission \gaia\ (\url{https://www.cosmos.esa.int/gaia}), processed by the \gaia\ Data Processing and Analysis Consortium (DPAC, \url{https://www.cosmos.esa.int/web/gaia/dpac/consortium}). Funding for the DPAC has been provided by national institutions, in particular the institutions participating in the \gaia\ Multilateral Agreement. This research has made use of the Exoplanet Follow-up Observation Program (ExoFOP; DOI: 10.26134/ExoFOP5) website and of the NASA Exoplanet Archive, both operated by the California Institute of Technology under contract with the National Aeronautics and Space Administration under the Exoplanet Exploration Program. The analysis code and the manuscript were developed with the assistance of an AI system (Claude, Anthropic). The author directed the work, verified the results and takes full responsibility for the content; every numerical result can be reproduced from the archived code and data \citep{DataRecord}.


\facilities{TESS, Gaia}

\software{astropy \citep{Astropy2022}, astrocut \citep{Brasseur2019}, batman \citep{Kreidberg2015}, emcee \citep{ForemanMackey2013}, lightkurve \citep{Lightkurve2018}, matplotlib \citep{Hunter2007}, numpy \citep{Harris2020}, scipy \citep{Virtanen2020}, tess-point \citep{Burke2020}, TRICERATOPS \citep{Giacalone2020}}

\section*{Data Availability}
The light curves of all transit windows, our photometry of Sectors~100--103, the measured transit times, the comparison with the QLP light curves of Sectors~100 and~102, the posterior summaries and chains, the next 30 predicted transits, and the analysis code are available on Zenodo \citep{DataRecord}. The pre-registered test of the 2027 transits is a separate Zenodo record \citep{PreReg}.

\bibliography{refs}{}
\bibliographystyle{aasjournal}

\end{document}
